\documentclass[10pt,a4paper]{amsart}
\usepackage[margin=19mm]{geometry}
\usepackage{amsmath,amssymb,mathtools}
\usepackage[scr=rsfs,cal=euler]{mathalfa}
\usepackage{xfrac}
\usepackage[unicode,hidelinks,bookmarks=false]{hyperref}
\newcommand{\ie}{i.e.\ }
\newcommand{\cf}{cf.\ }
\newcommand{\resp}{respectively\ }
\newcommand{\Subsection}{\subsection}
\providecommand{\nobreakdash}{\nobreak\hbox{-}}
\setlength{\emergencystretch}{2em}
\multlinegap0pt
\usepackage{amssymb}
\newtheorem{theorem}{Theorem}
\newtheorem{lemma}{Lemma}
\newcommand{\DD}{\mathrm D}
\newcommand{\HHH}{\mathscr{H}}
\DeclareMathOperator{\SL}{SL}
\def \ge {\geqslant}
\def \le {\leqslant}
\newcommand{\mumu}{\boldsymbol{\mu}}
\title{On $p$-Jordan constant of Cremona group of rank~$2$ in odd characteristic}
\author{Yifei Chen, Constantin Shramov}
\date{Source: arXiv:2509.15522v3 (CC BY 4.0)}
\begin{document}
\maketitle

\noindent\emph{Source and licence.} On $p$-Jordan constant of Cremona group of rank~$2$ in odd characteristic. Version arXiv:2509.15522v3 (CC BY 4.0), distributed by arXiv under the Creative Commons Attribution 4.0 International licence, http://creativecommons.org/licenses/by/4.0/, as recorded in the arXiv licence metadata for this exact version and shown on the abstract page https://arxiv.org/abs/2509.15522v3. Full licence notice: \texttt{licenses/arxiv-2509.15522-CC-BY-4.0.txt}.\par\smallskip

\noindent\textit{Editorial scope.} Counted unit: the article's lemma lemma:auxiliary-subgroups (source0.tex lines 754--800). The article's complete original proof of it is delivered with it (source0.tex lines 802--956, one \texttt{proof} environment of 155 lines). No mathematics is written, repaired or re-derived: the statement and the proof are the article's own lines, in the article's own notation, and no retained line is substituted or re-flowed.  Retained uncounted as the source-local context the proof needs: lines 741--746.  Nothing was omitted from any retained span, so the package carries no omission record.  No retained line contains a citation, so the wrapper carries no bibliography and no bibliography entry.  No retained line contains a figure environment, an image-inclusion command or drawing code, so no image is delivered and the package declares no asset dependency.  Editorial formatting only, in addition: an AMS \texttt{amsart} preamble that restates the article's own declarations and macros used by the retained lines, namely the environments \texttt{theorem}, \texttt{lemma} on separate counters, together with the standard packages those lines need.  The retained environments print in this document as Theorem 1, Lemma 1 in order of appearance, whereas the article prints the same items with the numbering of its own full document.  The complete native source of the article as distributed in the arXiv e-print for this exact version is bundled unchanged as \texttt{sources/185316/source0.tex}.
\begin{theorem}
\label{theorem:normal-in-a-p-group}
Let $p$ be a prime number, and let $G$ be a group of order~$p^n$.
Then~$G$ contains a normal abelian subgroup of order $p^m$ for some
positive integer~$m$ such that $m(m+1)\ge 2n$.
\end{theorem}

\begin{lemma}\label{lemma:auxiliary-subgroups}
Let $p$ be a prime number, and let $G$ be a finite group.
Set
\begin{equation}\label{eq:auxiliary-subgroups}
J=\begin{cases}
720, &\text{if\ } p\ge 7,\\
144, &\text{if\ } p=5,\\
10, &\text{if\ } p=3,\\
3, &\text{if\ } p=2.
\end{cases}
\end{equation}
The following assertions hold.
\begin{itemize}
\item[(i)] Suppose that $G\subset\mathfrak{S}_5$. Then $G$ contains a normal abelian subgroup
of order coprime to $p$ and index at most $J\cdot |G_{(p)}|^3$.
Moreover, one has~\mbox{$|G|\le J\cdot |G_{(p)}|^3$}, unless either $p=3$ and $G\cong\mumu_5\rtimes\mumu_4$, or $p=2$ and~\mbox{$G\cong\mumu_5$}.

\item[(ii)] Suppose that $G\subset \mumu_2^4\rtimes\mathfrak{S}_5$, where
the action of $\mathfrak{S}_5$ on $\mumu_2^4$ comes from the permutation representation.
Then $G$ contains a normal abelian subgroup
of order coprime to $p$ and index at most $J\cdot |G_{(p)}|^3$,
unless $p=3$ and~\mbox{$G\cong\mumu_2^4\rtimes (\mumu_5\rtimes\mumu_4)$}.

\item[(iii)] Suppose that $G\subset \mumu_2^4\rtimes\mathfrak{A}_5$, where
the action of $\mathfrak{A}_5$ on $\mumu_2^4$ comes from the permutation representation.
Then $G$ contains a normal abelian subgroup
of order coprime to $p$ and index at most $J\cdot |G_{(p)}|^3$.

\item[(iv)] Suppose that $G\subset\mathfrak{S}_6$.
Then $G$ contains a normal abelian subgroup
of order coprime to $p$ and index at most $J\cdot |G_{(p)}|^3$.

\item[(v)] Suppose that $G\subset\Gamma=\HHH_3\rtimes\SL_2(\mathbf{F}_3)$.
Then $G$ contains a normal abelian subgroup
of order coprime to $p$ and index at most $J\cdot |G_{(p)}|^3$,
unless $p=5$ and either $G=\Gamma$, or $|G|=162$.

\item[(vi)] Suppose that $G\subset\mumu_3^3\rtimes \mathfrak{S}_4$.
Then $G$ contains a normal abelian subgroup
of order coprime to $p$ and index at most $J\cdot |G_{(p)}|^3$.

\item[(vii)] Suppose that $G\subset\Gamma$, where $\Gamma$ is a group of order $576$.
Then $G$ contains a normal abelian subgroup
of order coprime to $p$ and index at most~\mbox{$J\cdot |G_{(p)}|^3$},
unless $p=5$ and $|G|\in\{192,288,576\}$.
\end{itemize}
\end{lemma}

\begin{proof}
Suppose that $G\subset \mathfrak{S}_5$. If $p\ge 5$, then $|G|<J$, so that
there is nothing to prove. Thus, we have to consider only the cases $p=2$ and $p=3$.
One has
$$
|G|\in\{1,2,3,4,5,6,8,10,12,20,24,60,120\}.
$$
Checking the orders of the $2$- and $3$-Sylow subgroups case by case, we see that the
inequality $|G|\le J\cdot |G_{(p)}|^3$ fails only if $p=3$ and $|G|=20$, and if
$p=2$ and~\mbox{$|G|=5$}. In the former case one has $G\cong\mumu_5\rtimes\mumu_4$,
and in the latter case~\mbox{$G\cong\mumu_5$}. This proves the second part of assertion~(i).
Now observe that if the bound~\mbox{$|G|\le J\cdot |G_{(p)}|^3$} holds, then the first part of assertion~(i)
holds as well. On the other hand, if $p=3$ and $G\cong\mumu_5\rtimes\mumu_4$, then
$G$ contains a normal abelian subgroup $\mumu_5$ of index $4<J$.
This completes the proof of assertion~(i).

Suppose that $G$ satisfies the assumptions of assertion (ii). Then $G$ fits into
an exact sequence
$$
1\to H\to G\to\bar{G}\to 1,
$$
where $H\subset \mumu_2^4$ and $\bar{G}\subset\mathfrak{S}_5$.
If $p\ge 3$, then it follows from assertion~(i) that
$$
\frac{|G|}{|H|}=|\bar{G}|\le J\cdot |\bar{G}_{(p)}|^3= J\cdot |G_{(p)}|^3,
$$
unless $p=3$ and $\bar{G}\cong\mumu_5\rtimes \mumu_4$. Observe that
$\mumu_2^4$ does not contain proper non-trivial $\mumu_5$-invariant subgroups.
This means that if $\bar{G}\cong\mumu_5\rtimes \mumu_4$, then we have either $G\cong\mumu_5\rtimes\mumu_4$, so that $G$ contains a normal subgroup $\mumu_5$
of index~\mbox{$4<J$}, or $G\cong\mumu_2^4\rtimes (\mumu_5\rtimes\mumu_4)$.
Thus, we are left with the case $p=2$.
In this case assertion~(i) implies that
$$
|G|=|H|\cdot |\bar{G}|=|H_{(2)}|\cdot |\bar{G}|\le |H_{(2)}|^3\cdot  J\cdot |\bar{G}_{(2)}|^3=
J\cdot |G_{(2)}|^3,
$$
unless $\bar{G}\cong \mumu_5$.
Since $\mumu_2^4$ does not contain proper non-trivial $\mumu_5$-invariant subgroups,
we see that in the latter case we have either $G\cong\mumu_5$, so that
$G$ is abelian and has order coprime to~$2$, or
$G\cong\mumu_2^4\rtimes \mumu_5$, so that
$$
|G|=40 < 3\cdot 2^{12}=J\cdot |G_{(2)}|^3.
$$
This proves assertion~(ii). Assertion~(iii) is a direct consequence of assertion~(ii).

Suppose that $G\subset\mathfrak{S}_6$.
If $p\ge 7$, then $|G|\le J$. If $p=5$ and $G$ is either $\mathfrak{S}_6$ or $\mathfrak{A}_6$,
then $|G|<J\cdot |G_{(5)}|^3$. If $G$ is neither $\mathfrak{S}_6$ nor $\mathfrak{A}_6$,
then $|G|\le 120$; in this case we have $|G|< J$ for $p=5$.
Thus, we have to consider only the cases $p=2$ and $p=3$.
One has
$$
|G|\in\{1,2,3,4,5,6,8,9,10,12,16,18,20,24,36,48,60,72,120,360,720\}.
$$
Checking the orders of the $2$- and $3$-Sylow subgroups case by case, we see that the
inequality $|G|\le J\cdot |G_{(p)}|^3$ fails only if $p=3$ and $|G|\in\{16,20\}$,
and if~\mbox{$p=2$} and $|G|\in\{5,9\}$. In the former case one has either $G\cong \mumu_2\times\DD_8$,
so that~$G$ contains a normal abelian subgroup $\mumu_2\times\mumu_4$ of index $2<J$,
or $G\cong \mumu_5\rtimes\mumu_4$, so that $G$ contains a normal abelian subgroup $\mumu_5$ of index $4<J$. In the latter case one has either $G\cong \mumu_5$, or $G\cong \mumu_3^2$;
in both cases $G$ is abelian and has order coprime to $2$.
This proves assertion~(iv).

Suppose that $G\subset \Gamma\cong\HHH_3\rtimes\SL_2(\mathbf{F}_3)$.
If $p\ge 7$, then $|G|<J$, so that there is nothing to prove.
The group $G$ fits into
an exact sequence
$$
1\to H\to G\to\bar{G}\to 1,
$$
where $H\subset \HHH_3$ and $\bar{G}\subset\SL_2(\mathbf{F}_3)$.
Suppose that $p=5$ and $G\neq \Gamma$; then~\mbox{$|G|\le 324$}.
If $H$ contains the center $Z\cong\mumu_3$ of $\HHH_3$, then $Z$ is a normal subgroup
of $G$, and its index in $G$ is
$$
\frac{|G|}{|Z|}\le \frac{324}{3}=108<J.
$$
If $H$ does not contain $Z$, then $H$ is abelian, and its index in $G$
is at most~\mbox{$|\bar{G}|=24<J$}. Thus, we have to consider only the cases $p=2$ and $p=3$.
One has
$$
|\bar{G}|\in\{1,2,3,4,6,8,24\}.
$$
If $p=3$, we conclude that $|\bar{G}|<J\cdot |\bar{G}_{(3)}|^3$, which means that
$$
|G|=|H|\cdot |\bar{G}|=|H_{(3)}|\cdot |\bar{G}|<|H_{(3)}|^3\cdot J\cdot |\bar{G}_{(3)}|^3=
J\cdot |G_{(3)}|^3.
$$

Thus, to prove assertion~(v) one can suppose that $p=2$. In this case we see that
$|\bar{G}|\le J\cdot |\bar{G}_{(2)}|^3$. Hence, if $H$ is abelian, then
we are done. Suppose that~$H$ is not abelian.
Then~\mbox{$H=\HHH_3$}, and $Z$ is a normal abelian subgroup of~$G$ of index
$$
9\cdot |\bar{G}|<J\cdot |\bar{G}_{(2)}|^3,
$$
unless $|\bar{G}|\in\{1,3,6\}$. If $|\bar{G}|=1$, then $G\cong \HHH_3$ contains a normal
abelian subgroup of index $3=J$. If $|\bar{G}|=3$, then $|G|=81$, and hence
$G$ contains a normal abelian subgroup of index $3$ by Theorem~\ref{theorem:normal-in-a-p-group}.
This completes the proof of assertion~(v).

Suppose that $G\subset\mumu_3^3\rtimes \mathfrak{S}_4$.
Then $G$ fits into
an exact sequence
$$
1\to H\to G\to\bar{G}\to 1,
$$
where $H\subset \mumu_3^3$ and $\bar{G}\subset\mathfrak{S}_4$.
If $p\neq 3$, then the index of the normal abelian subgroup $H$ in $G$
is at most
$$
|\bar{G}|\le J\cdot |\bar{G}_{(p)}|^3=J\cdot |G_{(p)}|^3
$$
by assertion~(i).
Thus, we have to consider only the case $p=3$.
One has
$$
|\bar{G}|\in\{1,2,3,4,6,8,12,24\}.
$$
In each of these cases we have $|\bar{G}|< J\cdot |\bar{G}_{(3)}|^3$,
so that
$$
|G|=|H|\cdot |\bar{G}|=|H_{(3)}|\cdot |\bar{G}|< |H_{(3)}|^3\cdot J\cdot |\bar{G}_{(3)}|^3=
J\cdot |G_{(3)}|^3.
$$
This proves assertion~(vi).

Suppose that $G\subset\Gamma$, where $\Gamma$ is a group of order $576$.
If $p\ge 7$, then $|G|<J$, so that there is nothing to prove.
If $p=5$ and the index of $G$ in $\Gamma$ is at least $4$, then $|G|\le 144=J$;
if the index of $G$ in $\Gamma$ is at most $3$, then
$$
|G|\in\{192,288,576\}.
$$
Thus, we have to consider only the cases $p=2$ and $p=3$.
Write $|G|=2^a3^b$, where $a\le 6$ and $b\le 2$. If
$p=3$, then
$$
|G|=10\cdot \frac{2^a}{10\cdot 3^{2b}}\cdot 3^{3b}=
J\cdot \frac{2^{a-1}}{5\cdot 3^{2b}}\cdot |G_{(3)}|^3\le J\cdot |G_{(3)}|^3,
$$
unless $2^{a-1}>5\cdot 3^{2b}$. The latter is possible only if $a\ge 4$ and $b=0$, which
gives~\mbox{$|G|\in\{16,32,64\}$}. In each of these cases $G$ contains a normal abelian subgroup of
order at least $8$ by Theorem~\ref{theorem:normal-in-a-p-group}, so that this
subgroup has index at most~\mbox{$8<J$} in $G$.
If $p=2$, then
$$
|G|=3\cdot \frac{3^{b-1}}{2^{2a}}\cdot 2^{3a}=
J\cdot \frac{3^{b-1}}{2^{2a}}\cdot |G_{(2)}|^3\le J\cdot |G_{(2)}|^3,
$$
unless $3^{b-1}>2^{2a}$. The latter is possible only if $b=2$ and $a\le 1$, which
gives~\mbox{$|G|\in\{9,18\}$}. In the former case $G$ is abelian, while in the latter
case the $3$-Sylow subgroup of $G$ is abelian, has index $2<J$ in $G$, and is normal
in~$G$. This proves assertion~(vii).
\end{proof}

\end{document}
